\documentclass[10pt]{article}

\usepackage[utf8]{inputenc}

\usepackage[T1]{fontenc}

\usepackage{amsmath}

\usepackage{amsfonts}

\usepackage{amssymb}

\usepackage[version=4]{mhchem}

\usepackage{stmaryrd}

\usepackage{bbold}

\usepackage{graphicx}

\usepackage[export]{adjustbox}

\graphicspath{ {./images/} }



\begin{document}

Значение параметра $\alpha$ (параметра регуляризации) должно быть согласовано с уровнем погрешности исходных данных. Его можно определить, напр., по невязке, т. е. из условия $\rho_{U}\left(A z_{\alpha}, \tilde{u}\right)=\delta$, если известно число $\delta$. Но возможны и др. способы определения $\alpha$ (см. [1]). Таким образом, параметр $\alpha$ должен зависеть от $\delta$ и $\tilde{u}$, $\alpha=\alpha(\delta, \tilde{u})$. Элемент $z_{\alpha(\delta, \tilde{u})}$ и принимается за приближенное решение уравнения $A z=\tilde{u}$. Это и есть одна из форм разработанного в [2], [3] Р. м. Аналогично строятся приближенные решения уравнений $\tilde{A} z=\tilde{u}$ с приближенно ваданным оператором $\tilde{A}$ и правой частью $\tilde{u}$. При этом минимизируется фуункцинал типа $M^{\alpha}[z, \tilde{A}$, u] (см., напр., [1]). Возможны и другие формы Р. м. и применение его к иным классам задач(см. [1]). Р. м. развит и для решения нелинейных задач (см. [1], [4]). Лum.: [1] Т и х о н о в А. Н., А р с е н и н В. Я., Методы решения некорректных задач, 2 изд., М., 1979; [2] т и х оНо в А. Н., "Докл. АН ССCP", 1963, т. 151, Nㅜ 3, с. 501-04; е в М. М., О некоторых некорректных задачах математической $\begin{array}{ll}\text { физики, Новосиб, } 1962 & \text { B. Я. Арсенин, А. Н. Тихонов. }\end{array}$



РЕГУ ЛЯРИЗАЦИЯ - использование той или иной формы отбора допустимых решений при построении устойчивых $\kappa$ исходной информации приближенных решений некорректно поставленных задач (см. также Некорректиые задачи и Регуляризации метод).



РЕГУЛЯРНАЯ ГРАНИЧНАЯ ТОЧКА - точка $y_{0}$ границы $\Gamma$ области $D$ евклидова пространства $\mathbb{R}^{n}$, $n \geqslant 2$, в к-рой для любой непрерывной на $\Gamma$ функции $f(y)$ обобщенное решение $u(x)$ Дирихле задачи в смысле Винера - Перрона (см. Перрона метод) принимает граничное значение $f\left(y_{0}\right)$, то есть



\[

\lim _{x \rightarrow y_{0}, x \in D} u(x)=f\left(y_{0}\right) .

\]



P. г. т. для области $D$ образуют множество $R$, в точках к-рого дополнение $C D=\mathbb{R} n \backslash D$ не является разреженным множеством; множество $\Gamma \backslash R$ ирреәулярных араничных точек есть полярное множество типа $F_{\sigma}$. Если все точки $\Gamma$ суть Р. г. т., то область $D$ наз. р ег у л я р н о й относительно задачи Дирихле.



Для того чтобы точка $y_{0} \in \Gamma$ была Р. Г. т., необходимо и достаточно, чтобы в пересечении $U_{0}=U \cap D$ области $D$ с нек-рой окрестностью $U$ точки $y_{0}$ существовал барьер, т. е. супергармонич. функция $\omega(x)>0$ в $U_{0}$ та-



\includegraphics[max width=\textwidth, center]{2023_07_08_f59f2bd261b8fd8756c4g-1}

б е г а). А. Лебег (H. Lebesgue, 1912) впервые показал, что при $n \geqslant 3$ вершина достаточно острого входящего в $D$ острия может не быть P. г. т.



Пусть



\[

E_{k}=\left\{x \in C D: 2^{-k} \leqslant\left|x-y_{0}\right| \leqslant 2^{-k+1}\right\}

\]



$c_{k}=C\left(E_{k}\right)$ - емкость $E_{k}$. Для того чтобы точка $y_{0} \in \Gamma$ была Р.г.т., необходимо и достаточно, чтобы расходился ряд



\[

\sum_{k=1}^{\infty} 2^{k(n-2)} c_{k}, n \geqslant 3

\]



или при $n=2$ ряд



\[

\sum_{k=1}^{\infty} 2^{k} c_{k}

\]



причем здесь



\[

E_{k}=\left\{x \in C D: 2^{k} \leqslant \ln \frac{1}{\left|x-y_{0}\right|} \leqslant 2^{k+1}\right\}

\]



(к р и т е р и й В и не р а).



При $n=2$ точка $y_{0} \in \Gamma$ является Р. г. т., если существует непрерывный путь $x(t), 0 \leqslant t \leqslant 1$, такой, что $x(1)=y_{0}, x(t) \in C D$ при $0 \leqslant t<1$. При $n \geqslant 3$ точка $y_{0} \in \Gamma$ является Р. г.т., если ее можно коснуться вершиной прямого кругового конуса, принадлежацего $C D$ в достаточно малой окрестности $y_{0}$. В случае области $D$ компактифицированного пространства $\overline{\mathbb{R}} n, n \geqslant 3$, бесконечно удаленная точка $\infty \in \Gamma$ всегда является Р. г. т.; при $n=2$ бесконечно удаленная точка $\infty \in \Gamma$ является P. г. т., если существует непрерывный путь $x(t), 0 \leqslant t<1$, такой, что $x(t) \in C D$ при $0 \leqslant t<1$ и $\lim _{t \rightarrow 1-0} x(t)=\infty$.



См. также Иррегулярная граничная точка.



Лит.: [1] К е Ј Д ы II М. В., "Успехи матем. наук", 1941, в. 8, с. 171-232; [2] Л а н Д к о Ф Н. С., Основы современной теории потенциала, М., 1966; [3] Х е и ма н У., К е н н ед и П., Субгармонические функции, пер. с англ., М. М., 1980.



РЕГУЛЯРНАЯ $\boldsymbol{p}$-ГРУППА - $p$-группа $G$ такая, что для любых ее элементов $a, b$ и любого целого $n=p^{\alpha}$ справедливо равенство



\[

(a b)^{n}=a^{n} b^{n} s_{1}^{n} \ldots s^{n},

\]



где $s_{1}, \ldots, s_{t}$ - нек-рые әлементы из ком.мутанта подгруппы, порожденной элементами $a$ и $b$. Подгруппы и факторгруппы Р. $p$-г. регулярны. Конеqная $p$-группа регулярна тогда и только тогда, когда для любых ее элементов $a$ и $b$ справедливо равенство



\[

a^{p} b^{p}=(a b)^{p} s^{p}

\]



где $s$ - нек-рый элемент коммутанта подгруппы, порожденной элементами $a$ и $b$.



Әлементы P. p-г. $G$, имеющие вид $a^{\alpha}, a \in G$, образуют характеристич. подгруппу $C^{\alpha}(G)$, а элементы порядка, не большего числа $p^{\alpha},-$ вполне характеристич. подгруппу $C_{\alpha}(G)$.



Примерами Р. p-г. являются любая $p$-группа, класс нильпотентности к-рой меньше $p$, а такжө любая $p$ группа порядка, не большего числа $p p$. Для любого $p$ существует нерегулярная $p$-группа порядка $p^{p+1}, \mathrm{a}$ именно, силовская подгруппа $S_{p}$ симметрич. группы $S\left(p^{2}\right)$ степени $p^{2}$ (она изоморфна сплетению циклич. группы порядка $p$ с самой собой).



Лит.: [1] Х о л л М., Теория групп, пер. с англ., М., 1962.



РЕГУЛЯРНАЯ МЕРА - мера, определенная на борелевской $\sigma$-алгебре $\mathfrak{Z}(T)$ топологич. пространства $T$ такая, что для любого борелевского множества $X \in \mathfrak{B}(T)$ и любого $\varepsilon>0$ найдется открытое множество $G \subset T$, покрывающее $X: X \subset G$ и $\mu(G X)<\varepsilon$. Равносильное определение: для любого $X \in \mathfrak{B}(T)$ и $\varepsilon>0$ найдется замкнутое множество $F \subset X$ такое, что $\underset{P}{\mu}(X$. Минлос.



РЕГУЛЯРНАЯ ОСОБАЯ ТОЧКА - понЯтіе теории обыкновенных линейных дифференциальных уравненийі с комплексным независимым переменным. Точка $a \in \mathbb{C}$ наз. Р. о. т. уравнения



\[

y^{(n)}+a_{1}(t) y^{(n-1)}+\ldots+a_{n}(t) y=0

\]



или системы



\[

\dot{z}=A(t) z, \quad z \in \mathbb{C}^{n},

\]



с аналитич, коэффициентами, если $a-$ изолированная особенность коэффициентов и все репения уравнения (1) или системы (2) растут не быстрее, чем $|t-a|^{d}$ для нек-рого $d \in \mathbb{R}$, когда $t$ стремится к $a$, оставаясь внутри произвольного острого угла с веріпиной $a$. Последнее ограничение вызвано тем, что в окрестности Р. о. т. решения являются неоднозначными аналитич. функциями и при $t \rightarrow a$ по произвольной кривой могут расти существенно быстрее, чем при стремлении $t \rightarrow a$ по лучу с вершиной $a$.



Для того чтобы особая точка коәффициентов уравнения (1) или системы (2) была Р. о. т., необходимо, чтобы она была полюсом, а не существенно особой точкой коэфффициентов. Для уравнений (1) имеет место у с л о в и е Ф у к с а: особая точка $t=0$ коәффициен-





\end{document}